% ==============================================================================
%  CHAPTER 3 — METHODOLOGY
%  Paste this block into your main document immediately before \begin{thebibliography}
% ==============================================================================

\section{Methodology}

This chapter presents the design and implementation of the proposed context-aware
privacy framework. The framework is structured as a five-stage sequential pipeline that
processes each frame of an egocentric video stream: face detection, multi-face tracking,
interaction cue extraction, context-aware classification, and selective anonymization.
Each stage is motivated by the design constraints established in the literature review:
CPU-only edge deployment, no biometric face recognition, and a privacy-safe
``bystander by default'' classification policy. The chapter concludes with a description
of the blanket-blurring baseline used for comparative evaluation and the metrics used to
assess performance.

\subsection{System Architecture}

The proposed system is organised as a modular, single-pass pipeline that processes
one video frame at a time. Figure~\ref{fig:pipeline} illustrates the data flow. A raw
BGR frame is passed sequentially through the detection, tracking, cue extraction,
classification, and anonymization stages; the output is an anonymized frame in which
only bystander faces are blurred, while interaction partner faces remain unmodified.

\begin{figure}[htbp]
  \centering
  \fbox{
    \begin{minipage}{0.9\linewidth}
      \centering
      \vspace{4pt}
      \texttt{Frame} $\rightarrow$
      \textbf{Detection} $\rightarrow$
      \textbf{Tracking} $\rightarrow$
      \textbf{Cue Extraction} $\rightarrow$
      \textbf{Classification} $\rightarrow$
      \textbf{Anonymization} $\rightarrow$
      \texttt{Output}
      \vspace{4pt}
    \end{minipage}
  }
  \caption{Overview of the five-stage egocentric privacy pipeline.}
  \label{fig:pipeline}
\end{figure}

A key architectural property is the separation of \emph{identity management} from
\emph{appearance analysis}. The tracking stage assigns and maintains stable numeric
identifiers for each face across frames; the cue extraction stage operates on the
current frame content given those identifiers. This separation means that no biometric
embedding or face recognition model is required at any stage. The numeric identifiers
are session-scoped and ephemeral, consistent with the Data Minimization principle
discussed in Section~2.6.

All components are implemented in Python 3.11 using OpenCV 4.x for frame I/O and
image operations, Ultralytics YOLOv8 for face detection, a custom SORT implementation
for tracking, and MediaPipe 0.10 for facial landmark extraction. The pipeline is
designed to execute on a single CPU core with no GPU or NPU dependency, targeting the
Raspberry Pi~5 (ARM Cortex-A76, 2.4~GHz) as the primary edge deployment hardware.

\subsection{Face Detection}

The first stage detects all faces present in each frame and produces axis-aligned
bounding boxes and detection confidence scores. Based on the empirical comparison
described in Section~2.4, YOLOv8n-face was selected as the detection backbone.
YOLOv8n-face is the nano variant of the YOLOv8 architecture fine-tuned on face
detection datasets, with a model weight file of approximately 6~MB. It operates as a
single-shot detector: the full frame is resized to a fixed input resolution of
$640 \times 640$ pixels, passed through the network, and non-maximum suppression (NMS)
is applied to the raw anchor outputs to produce the final set of detections.

A minimum confidence threshold of $\tau_{\mathrm{det}} = 0.45$ and an NMS IoU threshold
of $\tau_{\mathrm{NMS}} = 0.45$ were selected empirically to suppress spurious low-confidence
detections without discarding valid faces in partially occluded or poorly lit conditions.
A maximum of 20 simultaneous faces are tracked per frame, a limit that comfortably
exceeds the face count expected in typical egocentric interaction scenarios.

\subsection{Multi-Face Tracking}

The tracking stage associates the set of detections from the current frame with the set
of active tracks from the previous frame, maintaining a stable numeric \texttt{track\_id}
for each face across consecutive frames. Stable identities are essential for the
interaction cue extraction stage, which accumulates temporal evidence such as mouth
aspect ratio variance over a 30-frame window; without tracking, each frame's detections
would be anonymous and no per-person history could be maintained.

SORT (Simple Online and Realtime Tracking) \cite{sort2016} is used for this purpose. SORT
models each face's bounding box trajectory as a linear Kalman filter and uses the
Hungarian algorithm to solve the bipartite matching problem between the predicted
positions of existing tracks and the new detections. The IoU between predicted and
detected bounding boxes serves as the matching cost. Tracks that fail to receive a
matching detection are retained in a ``lost'' state for up to 90 frames (3~seconds at
30~fps) before being discarded, which prevents identity switches when a face is
temporarily occluded. A track is reported as confirmed only after receiving at least 3
consecutive detections, which prevents transient false detections from entering the
classification stage.

SORT is a motion-only tracker: it requires no face embedding, identity descriptor, or
appearance model. This is a critical property for the present system, as any
re-identification model that encodes facial appearance would constitute biometric
processing. The Kalman filter parameters use velocity-and-position state estimation with
a constant-velocity motion model, which is adequate for the relatively smooth,
continuous motion of faces in egocentric video at normal walking pace.

\subsection{Interaction Cue Extraction}

The interaction cue extraction stage computes four scalar features for each tracked face
that collectively characterise whether the face belongs to an active interaction partner
or a background bystander. The four cues are divided into two groups based on the
information they require.

\subsubsection*{Phase 2a Cues: Bounding-Box Geometry}

The first two cues are derived entirely from the bounding box produced by the detector
and require no additional computation. They exploit the center bias of egocentric video:
the camera wearer naturally orients their head toward their interlocutor, so an active
conversation partner's face tends to be large and centrally located in the frame.

\textbf{Proximity} (area ratio). Let $[x_1, y_1, x_2, y_2]$ denote the bounding box of
a tracked face, and let $W \times H$ be the frame dimensions. The face area ratio is
defined as:
\begin{equation}
  r_{\mathrm{area}} = \frac{(x_2 - x_1)(y_2 - y_1)}{W \cdot H}
  \label{eq:area_ratio}
\end{equation}
$r_{\mathrm{area}}$ captures how much of the frame the face occupies. A close conversation
partner produces values in the range $0.02$--$0.08$; a distant bystander produces
values below $0.003$. A threshold of $\tau_{\mathrm{area}} = 0.012$ separates these two
regimes, calibrated from the empirical distribution of per-track maximum area ratios
observed across the evaluation videos.

\textbf{Centrality} (center distance). The normalized distance of the face center from
the optical center of the frame is defined as:
\begin{equation}
  d_{\mathrm{center}} = \frac{\sqrt{(\bar{x} - W/2)^2 + (\bar{y} - H/2)^2}}
                             {\sqrt{(W/2)^2 + (H/2)^2}}
  \label{eq:center_dist}
\end{equation}
where $\bar{x} = (x_1 + x_2)/2$ and $\bar{y} = (y_1 + y_2)/2$ are the face center
coordinates. The denominator normalizes by the half-diagonal of the frame so that
$d_{\mathrm{center}} = 0$ at the frame center and $d_{\mathrm{center}} \approx 1$ at the
corners. A threshold of $\tau_{\mathrm{center}} = 0.45$ is used: faces within 45\% of
the frame center along the normalized diagonal are considered sufficiently central.

\subsubsection*{Phase 2b Cues: Facial Landmark Analysis}

The second two cues require the positions of specific facial landmarks, which are
extracted using MediaPipe Face Mesh. For each tracked face, the bounding box produced
by SORT is padded by 15\% on each side to ensure that the chin and forehead are
captured, and the resulting crop is passed to MediaPipe Face Mesh in
\texttt{static\_image\_mode} (i.e., each crop is treated independently, with no temporal
state maintained by Face Mesh itself). Face Mesh returns a set of 468 3D landmark
coordinates normalized to the crop dimensions. If Face Mesh fails to detect a face
within the crop, the Phase 2b cues are set to \texttt{None} and the classification
stage reverts to Phase 2a cues only for that track in that frame.

\textbf{Head-pose yaw} (orientation cue). Head-pose estimation is performed by solving
the Perspective-$n$-Point (PnP) problem \cite{pnp1995} using six key facial landmarks
whose positions in a canonical 3D face model are known: the nose tip, chin, left and
right eye outer corners, and left and right mouth corners (MediaPipe landmark indices
4, 152, 33, 263, 61, and 291 respectively). Let
$\mathbf{P}_i^{3D} \in \mathbb{R}^3$ denote the $i$-th 3D model point and
$\mathbf{p}_i^{2D} \in \mathbb{R}^2$ its corresponding image point in the crop.
The camera intrinsic matrix is approximated as:
\begin{equation}
  K = \begin{bmatrix} f & 0 & c_x \\ 0 & f & c_y \\ 0 & 0 & 1 \end{bmatrix},
  \quad f = W_{\mathrm{crop}}, \quad c_x = W_{\mathrm{crop}}/2, \quad c_y = H_{\mathrm{crop}}/2
  \label{eq:camera_matrix}
\end{equation}
where $W_{\mathrm{crop}} \times H_{\mathrm{crop}}$ are the crop dimensions and $f$ is
the approximate focal length set equal to the crop width. The iterative PnP solver
(\texttt{SOLVEPNP\_ITERATIVE}) minimizes the sum of squared reprojection errors:
\begin{equation}
  \min_{\mathbf{r}, \mathbf{t}} \sum_{i=1}^{6}
  \left\| \mathbf{p}_i^{2D} - \pi\!\left(K,\, \mathbf{r},\, \mathbf{t},\, \mathbf{P}_i^{3D}\right) \right\|^2
  \label{eq:pnp}
\end{equation}
where $\pi(\cdot)$ denotes the perspective projection. The resulting rotation vector
$\mathbf{r}$ is converted to a $3\times3$ rotation matrix $R$ via the Rodrigues formula,
and the RQ decomposition of $R$ yields the Euler angles $(\alpha, \beta, \gamma)$
representing pitch, yaw, and roll respectively. The yaw angle $\beta$ (rotation about
the vertical axis) is used as the orientation cue:
\begin{equation}
  \theta_{\mathrm{yaw}} = \beta \quad [\text{degrees}]
  \label{eq:yaw}
\end{equation}
A face with $|\theta_{\mathrm{yaw}}| > \tau_{\mathrm{yaw}} = 40^\circ$ is considered to
be looking away from the camera wearer and is therefore classified as a bystander
regardless of proximity or centrality (hard veto).

\textbf{Mouth Aspect Ratio} (V-VAD cue). The Mouth Aspect Ratio (MAR) quantifies the
degree to which the mouth is open and is used as a visual proxy for speech activity.
Using the inner lip landmarks from MediaPipe Face Mesh, three upper inner lip points
(indices 82, 13, 312) and three lower inner lip points (indices 87, 14, 317) are
averaged to obtain representative upper ($\bar{u}$) and lower ($\bar{l}$) lip
positions. The left and right inner lip corners are taken from indices 78 and 308. The
MAR is then:
\begin{equation}
  \mathrm{MAR} = \frac{\lVert \bar{u} - \bar{l} \rVert}
                      {\lVert p_{78} - p_{308} \rVert}
  \label{eq:mar}
\end{equation}
where $\lVert \cdot \rVert$ denotes the Euclidean norm. MAR approaches zero when the
mouth is closed; values above approximately 0.08 correspond to a mouth open sufficiently
to indicate speech.

\subsection{Context-Aware Classification}

The classification stage maps the four cues computed for a given tracked face in a given
frame to a binary label: \textsc{Interactor} (preserve) or \textsc{Bystander}
(anonymize). The design follows three principles: (i) \emph{bystander by default}, so
that uncertain cases are treated as bystanders to maximise privacy protection; (ii)
\emph{temporal hysteresis}, so that a single noisy frame cannot flip a face's label;
and (iii) \emph{landmark cues as veto and boost}, so that the orientation and V-VAD cues
can strengthen or override the geometric Phase 2a decision.

\subsubsection*{Per-Track State}

Because SORT provides stable \texttt{track\_id} values, the classifier maintains a
stateful record for each active track containing:
\begin{itemize}
  \item \texttt{interacting}: the current committed binary label (initially \textsc{Bystander});
  \item \texttt{enter\_streak}: the count of consecutive frames in which the instantaneous
        cues indicate \textsc{Interactor};
  \item \texttt{exit\_streak}: the count of consecutive frames in which the cues indicate
        \textsc{Bystander};
  \item \texttt{mar\_history}: a fixed-length deque of the $N_w = 30$ most recent MAR
        values for this track, forming the V-VAD rolling window.
\end{itemize}

\subsubsection*{Instantaneous Decision Function}

For each frame, the instantaneous decision function $f$ maps the current cue vector to a
tentative label. The function is evaluated in the following priority order:

\begin{enumerate}
  \item \textbf{Yaw veto.} If landmark data is available and
        $|\theta_{\mathrm{yaw}}| > 40^\circ$, return \textsc{Bystander} immediately.
        A face turned significantly away from the camera cannot be in direct face-to-face
        interaction with the wearer.

  \item \textbf{Proximity and centrality.} Evaluate the Phase 2a rule:
        \begin{equation}
          f_{\mathrm{geo}} = \bigl[r_{\mathrm{area}} \geq \tau_{\mathrm{area}}\bigr]
                             \wedge
                             \bigl[d_{\mathrm{center}} \leq \tau_{\mathrm{center}}\bigr]
          \label{eq:geo_rule}
        \end{equation}
        If $f_{\mathrm{geo}}$ is true, return \textsc{Interactor} (subject to V-VAD
        boost below).

  \item \textbf{V-VAD boost.} The V-VAD score is computed as the fraction of frames in
        the rolling window with $\mathrm{MAR} \geq \tau_{\mathrm{MAR}} = 0.08$:
        \begin{equation}
          s_{\mathrm{vad}} = \frac{1}{|\mathcal{W}|} \sum_{m \in \mathcal{W}}
                             \mathbf{1}\!\left[m \geq \tau_{\mathrm{MAR}}\right]
          \label{eq:vvad}
        \end{equation}
        where $\mathcal{W}$ is the set of MAR values currently in the history deque.
        If $s_{\mathrm{vad}} \geq \tau_{\mathrm{vad}} = 0.30$ and
        $r_{\mathrm{area}} \geq \tau_{\mathrm{area}}$, the face is classified as
        \textsc{Interactor} even if $d_{\mathrm{center}} > \tau_{\mathrm{center}}$
        (with the centrality threshold relaxed to $1.2 \cdot \tau_{\mathrm{center}}$).
        This handles the case where an active conversation partner has drifted slightly
        off the optical center during animated speech.

  \item \textbf{Default.} If none of the above conditions is satisfied, return
        \textsc{Bystander}.
\end{enumerate}

\subsubsection*{Temporal Hysteresis and Latching}

The instantaneous label returned by $f$ updates the enter or exit streak for the track.
A label transition from \textsc{Bystander} to \textsc{Interactor} is committed only
when the enter streak reaches $\tau_{\mathrm{enter}} = 2$ consecutive frames. A special
case applies: if $r_{\mathrm{area}} \geq 0.04$ (a face covering at least 4\% of the
frame), the transition is committed immediately regardless of the streak, since such
proximity is an unambiguous indicator of direct face-to-face engagement.

Once a track has been committed as \textsc{Interactor}, the label is \emph{latched}: it
is not reverted to \textsc{Bystander} even if the instantaneous cues subsequently fail.
This design choice prevents flickering between labels when a conversation partner
momentarily looks away or is partially occluded, and is consistent with the observation
that once a direct interaction has been established the wearer's engagement partner
remains socially present throughout the interaction. Latching represents a deliberate
trade-off in favour of interaction preservation rate (IPR) over privacy success rate
(PSR), reflecting the system's utility objective.

\subsection{Selective Anonymization}

For each frame, all faces classified as \textsc{Bystander} are anonymized by applying
a Gaussian blur to the face region. The bounding box is first expanded by a padding
factor of 10\% on each side to ensure that hair and ears are covered in addition to the
central face region. A Gaussian kernel of size $51 \times 51$ pixels with a
correspondingly large standard deviation is applied to the padded region, which
produces a smooth, visually consistent blur that eliminates identifiable facial features
while preserving the presence and approximate shape of the person in the scene. Faces
classified as \textsc{Interactor} are passed to the output frame unmodified.

The selective output is produced alongside two reference streams: a \emph{blanket-blur}
stream in which all detected faces are blurred regardless of classification (the
baseline), and a \emph{debug} stream in which bounding boxes are annotated in green
(\textsc{Interactor}) or red (\textsc{Bystander}) to facilitate visual inspection of
the classification decisions.

\subsection{Baseline: Blanket-Blurring}

The baseline system applies identical Gaussian blur to every detected face in every frame,
making no use of the interaction cues. This corresponds to the current state-of-the-art
approach used in applications such as Google Street View and EgoBlur. The baseline uses
the same YOLOv8n-face detector and SORT tracker as the proposed system, ensuring that any
performance difference is attributable solely to the selective classification layer rather
than to differences in the underlying detection quality.

The proposed system is designed to outperform the baseline on two dimensions: the
Interaction Preservation Rate, which the baseline degrades to zero (all partners are
blurred), and the per-frame computational load, since selective anonymization skips the
Gaussian convolution for faces classified as \textsc{Interactor}.

\subsection{Evaluation Metrics}

The performance of the proposed system is evaluated using two primary task-level metrics,
supplemented by computational efficiency metrics.

\textbf{Privacy Success Rate (PSR)} measures the proportion of bystander faces that
are correctly anonymized:
\begin{equation}
  \mathrm{PSR} = \frac{TP_B}{TP_B + FN_B}
  \label{eq:psr}
\end{equation}
where $TP_B$ is the number of bystander face-frames correctly blurred and $FN_B$ is the
number of bystander face-frames incorrectly left unblurred. PSR captures the system's
ability to protect bystander privacy.

\textbf{Interaction Preservation Rate (IPR)} measures the proportion of interaction
partner faces that are correctly left unblurred:
\begin{equation}
  \mathrm{IPR} = \frac{TP_I}{TP_I + FN_I}
  \label{eq:ipr}
\end{equation}
where $TP_I$ is the number of partner face-frames correctly preserved and $FN_I$ is the
number of partner face-frames incorrectly blurred. IPR captures the system's utility
preservation.

Ground-truth labels (bystander or interaction partner, per face per frame) are
established by manual annotation of the evaluation video clips using a purpose-built
annotation tool. Both metrics are reported across all evaluation clips and separately for
the bystander-only and interaction video subsets, enabling analysis of false-positive
(over-anonymization) and false-negative (under-anonymization) failure modes.

\textbf{Computational efficiency} is reported as average frames per second (FPS) over
the full evaluation set, and as per-stage latency (in milliseconds) measured using
high-resolution wall-clock timing (\texttt{time.perf\_counter}). These measurements are
conducted on a desktop CPU to establish a baseline, with Raspberry Pi~5 performance
estimated by applying a hardware scaling factor derived from published ARM
Cortex-A76 inference benchmarks.

Table~\ref{tab:params} summarises the key threshold parameters used in the classification
stage.

\begin{table}[htbp]
\begin{singlespace}
\centering
\caption{Classification stage threshold parameters and their values.}
\label{tab:params}
\renewcommand{\arraystretch}{1.3}
\begin{tabular}{@{}llll@{}}
\toprule
\textbf{Parameter} & \textbf{Symbol} & \textbf{Value} & \textbf{Cue} \\
\midrule
Minimum area ratio       & $\tau_{\mathrm{area}}$    & $0.012$   & Proximity \\
Maximum center distance  & $\tau_{\mathrm{center}}$  & $0.45$    & Centrality \\
Instant-partner area     & --                         & $0.040$   & Proximity \\
Enter streak             & $\tau_{\mathrm{enter}}$   & $2$ frames & Hysteresis \\
Maximum yaw              & $\tau_{\mathrm{yaw}}$     & $40^\circ$ & Orientation \\
MAR threshold            & $\tau_{\mathrm{MAR}}$     & $0.08$    & V-VAD \\
V-VAD window             & $N_w$                     & $30$ frames & V-VAD \\
V-VAD active ratio       & $\tau_{\mathrm{vad}}$     & $0.30$    & V-VAD \\
\bottomrule
\end{tabular}
\end{singlespace}
\end{table}

% ── Add these to \begin{thebibliography} in your main document ─────────────────

% \bibitem{sort2016}
% A. Bewley, Z. Ge, L. Ott, F. Ramos, and B. Upcroft, ``Simple online and realtime
% tracking,'' in \emph{2016 IEEE International Conference on Image Processing (ICIP)},
% pp.~3464--3468, IEEE, 2016.
%
% \bibitem{pnp1995}
% L. Quan and Z. Lan, ``Linear N-point camera pose determination,''
% \emph{IEEE Transactions on Pattern Analysis and Machine Intelligence},
% vol.~21, no.~8, pp.~774--780, 1999.
